% status: SKELETON
% Chapter 2 of the second edition. Plan: docs/STRUCTURE.md. Rules: AUTHORING.md.
\chapter{The Forward Pass, Compressed}\label{ch:forward-pass}

\begin{ChapterIntent}
Before we can price a technique we need to know what one decode step
computes, in the terms an engine sees: operators, shapes, bytes and state.
This chapter walks a token through a modern decoder --- embedding,
normalization, attention over the cache, the feed-forward network or its
mixture of experts, the output projection and sampling --- and counts, for
each operator, the parameters it reads, the FLOPs it performs and the state it
touches. The result is the two rules the rest of the book leans on: a forward
pass costs about two FLOPs per active parameter per token, and batch-one decode
reads about every weight once per token. Appendix~A builds the same
computation from scratch.
\end{ChapterIntent}

\OpeningSection{The bytes a token actually touches}

\NapkinSection{Napkin math: parameters, FLOPs and bytes per layer}

\section{Tokens in, logits out}

\section{One decoder layer}

\section{Attention during decode: one query, many keys}

\section{The feed-forward network and its gate}

\section{Mixture of experts: a router in front of the FFN}

\section{Normalization, residuals and position}

\section{Logits, sampling and the loop}

\section{A shape-annotated map of one step}

\LedgerSection

\LabSection{count a real model from its file}

\HappenedSection{the embedding table that is read one row at a time}

\FrontierSection{}

\ExercisesSection
